\documentclass[12pt]{article}
\usepackage[utf8]{inputenc}
\usepackage{amsmath, amssymb, geometry, graphicx, setspace}
\geometry{letterpaper, margin=1in}
\setstretch{1.15}

\title{\textbf{Interplanetary Trajectories: \\ A Step-by-Step Guide to the Patched Conic Method}}
\author{Module: Astrodynamics}
\date{}

\begin{document}

\maketitle

\section*{The Challenge of Interplanetary Flight}
An interplanetary mission is fundamentally an $N$-body problem. A spacecraft traveling from Earth to Mars is gravitationally influenced simultaneously by the Earth, the Sun, Mars, and other celestial bodies. Because there is no general analytical solution to the $N$-body problem, we rely on a powerful simplification known as the \textbf{Patched Conic Approximation}.

This method relies on the fact that gravity obeys an inverse-square law. Close to a planet, the planet's gravity completely dominates the spacecraft's motion. Far from a planet, the Sun's gravity dominates. By dividing the solar system into distinct gravitational zones, we break the complex journey into three solvable $2$-body problems (conic sections) and ``patch'' them together at their boundaries.

\vspace{0.5cm}

\subsection*{The Sphere of Influence (SOI)}
The boundary where we switch our primary frame of reference from the planet to the Sun (or vice versa) is called the \textit{Sphere of Influence} ($r_{SOI}$). At this distance, the perturbing effect of the Sun on the planet-centric orbit roughly equals the perturbing effect of the planet on the Sun-centric orbit. 
\begin{equation}
    r_{SOI} = R_p \left( \frac{m_{planet}}{m_{sun}} \right)^{2/5}
\end{equation}
where $R_p$ is the distance from the Sun to the planet. While the SOI is vast from a human perspective (Earth's SOI is about 925,000 km), it is a mere speck on the scale of the solar system. Therefore, when looking at the Sun-centric (heliocentric) leg of the journey, the SOI is treated simply as a point coinciding with the planet's center.

\newpage

% ---------------------------------------------------------
% STEP 1
% ---------------------------------------------------------
\section*{Step 1: The Heliocentric Transfer (Phase 2)}
\textit{Counterintuitively, we must design the middle of the journey first to know how fast we need to leave the departure planet.}

In this phase, we ignore the planets' gravity and treat them as moving points. The spacecraft travels on a heliocentric transfer ellipse from the orbit of Planet 1 ($R_1$) to the orbit of Planet 2 ($R_2$). 

For a highly efficient Hohmann transfer, the semi-major axis of the transfer ellipse is $a_{tx} = (R_1 + R_2)/2$. Using the vis-viva equation, we calculate the absolute velocity the spacecraft must have with respect to the Sun ($V_{tx}$) at the moment it departs Planet 1:
\begin{equation}
    V_{tx} = \sqrt{\mu_{sun} \left( \frac{2}{R_1} - \frac{1}{a_{tx}} \right)}
\end{equation}

However, Planet 1 is already moving around the Sun with velocity $V_1 = \sqrt{\mu_{sun}/R_1}$. The difference between the required transfer speed and the planet's existing speed is the \textbf{Hyperbolic Excess Velocity ($v_\infty$)}. This is the relative speed the spacecraft must possess as it crosses the boundary of Planet 1's Sphere of Influence:
\begin{equation}
    v_{\infty} = | V_{tx} - V_1 |
\end{equation}
\textit{Note: We use uppercase $V$ for heliocentric (absolute) velocities, and lowercase $v$ for planetocentric (relative) velocities.}


% ---------------------------------------------------------
% STEP 2
% ---------------------------------------------------------
\section*{Step 2: Planetocentric Departure (Phase 1)}
Now we zoom in on Planet 1. The $v_\infty$ calculated above becomes our target boundary condition. The spacecraft must exit the SOI with a residual speed of $v_\infty$. To achieve this, it must travel on a \textbf{departure hyperbola} relative to Planet 1.

Assuming the spacecraft starts in a circular parking orbit of radius $r_p$ (where $r_p = R_{planet} + altitude$), its initial parking velocity is:
\begin{equation}
    v_c = \sqrt{\frac{\mu_1}{r_p}}
\end{equation}

By conservation of energy, the specific energy of the spacecraft at the parking orbit must equal the specific energy at the edge of the SOI (infinity). 
\begin{equation}
    \frac{v_p^2}{2} - \frac{\mu_1}{r_p} = \frac{v_{\infty}^2}{2} - 0
\end{equation}

Solving for $v_p$, the required velocity at the periapsis of the hyperbola (the burnout velocity), we get:
\begin{equation}
    v_p = \sqrt{v_{\infty}^2 + \frac{2\mu_1}{r_p}}
\end{equation}

The engine burn ($\Delta V_{depart}$) required to place the spacecraft onto this escape trajectory is simply the difference between the required burnout velocity and the velocity the spacecraft already has in its parking orbit:
\begin{equation}
    \Delta V_{depart} = v_p - v_c = \sqrt{v_{\infty}^2 + \frac{2\mu_1}{r_p}} - \sqrt{\frac{\mu_1}{r_p}}
\end{equation}


% ---------------------------------------------------------
% STEP 3
% ---------------------------------------------------------
\section*{Step 3: Planetocentric Arrival (Phase 3)}
The final phase is the mirror image of departure. The spacecraft arrives at Planet 2's Sphere of Influence. Relative to the Sun, its arrival velocity is $V_{arr}$. Planet 2 is moving around the Sun with velocity $V_2 = \sqrt{\mu_{sun}/R_2}$. 

The relative arrival velocity (the hyperbolic excess speed at Planet 2) is:
\begin{equation}
    v_{\infty, arr} = | V_{arr} - V_2 |
\end{equation}

As the spacecraft falls into Planet 2's gravity well, it accelerates along an \textbf{arrival hyperbola}. If the mission requires capturing into a circular parking orbit of radius $r_{cap}$ around Planet 2, we must apply the brakes at the periapsis of the hyperbola. 

First, we calculate the speed the spacecraft will naturally achieve at periapsis due to gravity falling in from infinity:
\begin{equation}
    v_{p, arr} = \sqrt{v_{\infty, arr}^2 + \frac{2\mu_2}{r_{cap}}}
\end{equation}

We only want to be traveling at the circular capture speed:
\begin{equation}
    v_{c, arr} = \sqrt{\frac{\mu_2}{r_{cap}}}
\end{equation}

Therefore, the retro-burn required to slow down and capture into the target orbit is:
\begin{equation}
    \Delta V_{capture} = v_{p, arr} - v_{c, arr}
\end{equation}

\vspace{0.5cm}

\section*{Mission Summary}
By breaking the interplanetary trajectory into these three segments, the total propellant budget for the mission can be accurately estimated by summing the engine burns:
\begin{equation}
    \Delta V_{total} = \Delta V_{depart} + \Delta V_{capture}
\end{equation}

This Patched Conic Approximation forms the baseline for all preliminary interplanetary mission design and allows engineers to evaluate launch windows, fuel requirements, and vehicle mass fractions before moving to numerical $N$-body simulations.

\end{document}